\chapter{Conclusion and Future Work}
\label{chap:conclusion}

% ============================================================================
% CHAPTER PURPOSE
% Close the argument. Do not introduce new evidence, equations or literature.
% Keep this chapter comparatively short and decisive.
% ============================================================================

\section{Research Aim and Approach}
\label{sec:conclusion-aim}

% Restate the aim in past tense and summarise the approach in 1--2 paragraphs:
% transparent deterministic model, software implementation, historic validation
% and sensitivity analysis. Do not repeat chapter-by-chapter contents.

\section{Principal Findings}
\label{sec:principal-findings}

% VERY HIGH PRIORITY. State only findings supported by Chapter 6, ideally in
% order of the research questions:
% - level of agreement with historic cases;
% - dominant parameters/mechanisms;
% - conditions under which predictions became unreliable;
% - value/limitations of race trace interpretation.
% Include numbers where they are the clearest summary, but no new tables.

\section{Achievement of Objectives}
\label{sec:objective-achievement-conclusion}

% Briefly revisit each numbered objective from Section
% \ref{sec:research-objectives}. Avoid duplicating the detailed Chapter 6 table;
% synthesise what was fully achieved, partially achieved or excluded.

\section{Contributions}
\label{sec:contributions}

% Be accurate and proportionate; "contribution" does not require claiming a new
% universal theory.
%
% Methodological/engineering contributions may include:
% - transparent decomposition of undercut gain into auditable components;
% - explicit treatment of common pit-loss cancellation;
% - progressive minimum-to-extended model and evaluation method;
% - evidence about parameter sensitivity for the selected cases.
%
% Practical contributions may include:
% - working historic-analysis application;
% - annotated race trace linking prediction and actual outcome;
% - reusable parameterised workflow.
% Only claim items actually delivered and evaluated.

\section{Recommendations for Use}
\label{sec:recommendations}

% MEDIUM PRIORITY but valuable.
% State the operating conditions in which the model is most defensible and the
% checks a user should perform before acting on its result. Emphasise decision
% support and safety margin under uncertain inputs.

\section{Future Work}
\label{sec:future-work}

% Prioritise extensions by the limitation they address and expected value.
% Suggested order:
% 1. improve degradation estimation/non-linear or stint-specific fitting;
% 2. quantify uncertainty and produce prediction ranges;
% 3. dynamic rejoin/traffic and multi-car interaction;
% 4. model overcut explicitly using the same comparative framework;
% 5. Safety Car/VSC/weather/probabilistic whole-race simulation;
% 6. live timing integration and prospective validation.
% For each: limitation -> proposed method -> expected benefit.
% Avoid a disconnected wish list (especially "use AI/ML" without justification).

\section{Final Statement}
\label{sec:final-statement}

% One concise closing paragraph that answers the original aim and leaves a
% balanced statement of usefulness and limits.
